\sectionheading{5.9 Python}

Python একটি \textbf{High Level, General Purpose Programming Language}।
এর Syntax তুলনামূলক সহজ, পরিষ্কার ও পাঠযোগ্য হওয়ায় নতুন
Programmer-দের জন্য এটি শেখা সহজ।

একই সঙ্গে Web Development, Automation, Data Analysis,
Artificial Intelligence ইত্যাদি আধুনিক ক্ষেত্রেও Python
ব্যাপকভাবে ব্যবহৃত হয়।

%================================================
% Definition
%================================================
\booksubheading{Python কাকে বলে?}

\begin{definitionbox}

\textbf{\color{darkTeal}সংজ্ঞা}

Python হলো একটি \textbf{High Level, General Purpose এবং
Interpreter-ভিত্তিক Programming Language}, যা সহজ ও পাঠযোগ্য
Syntax-এর জন্য পরিচিত।

\end{definitionbox}

\textbf{সহজভাবে মনে রাখবে—}

\begin{center}

\textbf{
Python
=
High Level
+
General Purpose
+
সহজ Syntax
}

\end{center}

%================================================
% History
%================================================
\booksubheading{Python-এর সংক্ষিপ্ত ইতিহাস}

Python সম্পর্কে কয়েকটি তথ্য MCQ-এর জন্য গুরুত্বপূর্ণ—

\begin{compacttwocolpoints}

\item Python-এর স্রষ্টা \textbf{Guido van Rossum}।

\item Python-এর প্রথম প্রকাশ ঘটে \textbf{1991 সালে}।

\item Python একটি \textbf{High Level Programming Language}।

\item এটি একটি \textbf{General Purpose Programming Language}।

\item Python-এর নামটি সরাসরি সাপের নাম থেকে নেওয়া হয়নি;
নামকরণে \textbf{Monty Python}-এর প্রভাব ছিল।

\end{compacttwocolpoints}

\begin{rememberbox}

\textbf{\color{deepOrange}মনে রাখবে}

\begin{center}

\textbf{
Guido van Rossum
\(\rightarrow\)
Python
\(\rightarrow\)
1991
}

\end{center}

\end{rememberbox}

%================================================
% Simple Example
%================================================
\booksubheading{Python-এর একটি সহজ উদাহরণ}

Screen-এ \textenglish{\ttfamily Hello} প্রদর্শনের জন্য
Python-এ লেখা যায়—

\begin{codebox}

\begin{center}
{\textenglish{\ttfamily print("Hello")}}
\end{center}

\end{codebox}

\textbf{Output:}

\begin{codebox}

\begin{center}
{\textenglish{\ttfamily Hello}}
\end{center}

\end{codebox}

এই ছোট Program থেকেই বোঝা যায় Python-এর Syntax তুলনামূলক
সহজ ও সংক্ষিপ্ত।

%================================================
% Basic Features
%================================================
\booksubheading{Python-এর প্রাথমিক বৈশিষ্ট্য}

\begin{tabularx}{\linewidth}{|>{\bfseries}p{0.27\linewidth}|Y|}
\hline

\rowcolor{softBlue}
\multicolumn{1}{|c|}{\textbf{বৈশিষ্ট্য}}
&
\multicolumn{1}{c|}{\textbf{সহজ ধারণা}}
\\
\hline

High Level
&
মানুষের জন্য তুলনামূলক সহজে লেখা ও বোঝা যায়।
\\
\hline

General Purpose
&
বিভিন্ন ধরনের Problem ও Application-এ ব্যবহার করা যায়।
\\
\hline

Interpreter-based
&
সাধারণত Interpreter-এর সাহায্যে Program চালানো হয়।
\\
\hline

Simple Syntax
&
Code তুলনামূলক সহজ, পরিষ্কার ও পাঠযোগ্য।
\\
\hline

Portable
&
বিভিন্ন Operating System-এ ব্যবহার করা যায়।
\\
\hline

Open Source
&
বিনামূল্যে ব্যবহার ও উন্নয়নের সুযোগ রয়েছে।
\\
\hline

Large Library
&
অনেক Built-in Module, Library ও Package পাওয়া যায়।
\\
\hline

\end{tabularx}

Python-এর বৈশিষ্ট্য পরবর্তী নির্দিষ্ট Topic-এ আরও বিস্তারিত
আলোচনা করা হবে।

%================================================
% Interpreter
%================================================
\booksubheading{Python-এ Interpreter-এর ভূমিকা}

Python Program সাধারণত Interpreter-এর সাহায্যে Execute করা হয়।

সহজভাবে—

\begin{center}
\begin{adjustbox}{max width=0.98\linewidth}
\fcolorbox{darkBlue}{softBlue}{
    \parbox[c][8mm][c]{0.24\linewidth}{
        \centering\textbf{Python\\Source Code}
    }
}
\quad
\(\longrightarrow\)
\quad
\fcolorbox{mainTeal}{softGreen}{
    \parbox[c][8mm][c]{0.24\linewidth}{
        \centering\textbf{Python\\Interpreter}
    }
}
\quad
\(\longrightarrow\)
\quad
\fcolorbox{darkBlue}{softBlue}{
    \parbox[c][8mm][c]{0.20\linewidth}{
        \centering\textbf{Execution}
    }
}
\end{adjustbox}
\end{center}

Interpreter Python Source Code-এর Instruction বুঝে Program
Execute করতে সাহায্য করে।

\begin{rememberbox}

\textbf{\color{deepOrange}মনে রাখবে}

Python-কে সাধারণভাবে একটি
\textbf{Interpreted Language} বলা হয়।

Interpreter সম্পর্কে পরবর্তী Translator Program অংশে আরও
বিস্তারিত আলোচনা করা হবে।

\end{rememberbox}

%================================================
% Uses
%================================================
\booksubheading{Python কোথায় ব্যবহার করা হয়?}

Python-এর ব্যবহার অনেক বিস্তৃত। যেমন—

\begin{compacttwocolpoints}

\item সাধারণ Programming ও Problem Solving-এ।

\item Web Development-এ।

\item Automation ও Scripting-এ।

\item Data Analysis-এ।

\item Scientific Computing-এ।

\item Machine Learning-এ।

\item Artificial Intelligence-এ।

\item Data Science-এ।

\item শিক্ষা ও গবেষণায়।

\end{compacttwocolpoints}

%================================================
% Why Popular
%================================================
\booksubheading{Python জনপ্রিয় হওয়ার কারণ}

Python জনপ্রিয় হওয়ার কয়েকটি প্রধান কারণ—

\begin{compacttwocolpoints}

\item Syntax সহজ ও পরিষ্কার।

\item কম Code-এ অনেক কাজ করা যায়।

\item নতুনদের জন্য শেখা তুলনামূলক সহজ।

\item অনেক Library ও Package পাওয়া যায়।

\item বিভিন্ন Operating System-এ ব্যবহার করা যায়।

\item Web, Automation, Data Science এবং AI-সহ বিভিন্ন ক্ষেত্রে
ব্যবহার করা যায়।

\end{compacttwocolpoints}

%================================================
% Advantages
%================================================
\booksubheading{Python-এর সুবিধা}

\begin{compacttwocolpoints}

\item Program লেখা ও পড়া সহজ।

\item Development তুলনামূলক দ্রুত করা যায়।

\item Debugging সহজ।

\item Code সাধারণত সংক্ষিপ্ত ও পাঠযোগ্য হয়।

\item Library ও Framework-এর সংখ্যা অনেক।

\item বিভিন্ন ধরনের Application তৈরিতে ব্যবহার করা যায়।

\item নতুন Programmer-দের জন্য শেখা সুবিধাজনক।

\end{compacttwocolpoints}

%================================================
% Limitations
%================================================
\booksubheading{Python-এর কিছু সীমাবদ্ধতা}

\begin{itemize}

\item অনেক ক্ষেত্রে Compiled Low Level Language-এর তুলনায়
Execution ধীর হতে পারে।

\item Hardware-এর ওপর Low Level Language-এর মতো সরাসরি
নিয়ন্ত্রণ কম।

\item কিছু Resource-sensitive Application-এর জন্য সবসময়
সেরা পছন্দ নাও হতে পারে।

\end{itemize}

%================================================
% Java vs Python
%================================================
\booksubheading{Java ও Python-এর প্রাথমিক পার্থক্য}

\begin{tabularx}{\linewidth}{|>{\bfseries}p{0.23\linewidth}|Y|Y|}
\hline

\rowcolor{softBlue}
\multicolumn{1}{|c|}{\textbf{বিষয়}}
&
\multicolumn{1}{c|}{\textbf{Java}}
&
\multicolumn{1}{c|}{\textbf{Python}}
\\
\hline

ভাষার স্তর
&
High Level
&
High Level
\\
\hline

ধরন
&
General Purpose
&
General Purpose
\\
\hline

প্রধান বৈশিষ্ট্য
&
Object-Oriented, JVM ও Bytecode
&
সহজ Syntax ও Interpreter-based ব্যবহার
\\
\hline

Syntax
&
তুলনামূলক বেশি নিয়মবদ্ধ
&
তুলনামূলক সহজ ও সংক্ষিপ্ত
\\
\hline

স্রষ্টা
&
James Gosling
&
Guido van Rossum
\\
\hline

প্রথম প্রকাশ
&
1995
&
1991
\\
\hline

\end{tabularx}

%================================================
% Machine Independence
%================================================
\booksubheading{Python কি Machine Independent?}

Python Source Code কোনো একটি নির্দিষ্ট Processor-এর Binary
Instruction দিয়ে লেখা হয় না।

প্রয়োজনীয় Python Environment বা Interpreter থাকলে একই
Source Code বিভিন্ন Operating System-এ ব্যবহার করা সম্ভব হতে পারে।

তাই Python-এর \textbf{Portability তুলনামূলক বেশি}।

\begin{rememberbox}

\textbf{\color{deepOrange}খেয়াল করো}

Machine Independent বা Portable বলার অর্থ এই নয় যে কোনো
Environment ছাড়াই Python Program সব Computer-এ চলবে।

Program চালানোর জন্য প্রয়োজনীয়
\textbf{Python Interpreter/Environment} থাকতে হবে।

\end{rememberbox}

%================================================
% Final Reminder
%================================================
\begin{rememberbox}

\textbf{\color{deepOrange}মনে রাখবে}

\begin{itemize}

\item Python একটি \textbf{High Level Programming Language}।

\item এটি একটি \textbf{General Purpose Language}।

\item Python-এর স্রষ্টা \textbf{Guido van Rossum}।

\item Python প্রথম প্রকাশিত হয় \textbf{1991 সালে}।

\item Python-কে সাধারণভাবে
\textbf{Interpreter-based Language} বলা হয়।

\item Python-এর Syntax সহজ, পরিষ্কার ও পাঠযোগ্য।

\item Python-এর নামের সঙ্গে \textbf{Monty Python}-এর সম্পর্ক রয়েছে।

\item Python Web Development, Automation, Data Science,
Machine Learning ও Artificial Intelligence-এ ব্যবহৃত হয়।

\item Python বিভিন্ন Platform-এ ব্যবহারযোগ্য এবং এর
Portability তুলনামূলক বেশি।

\item Python-এর বিস্তারিত বৈশিষ্ট্য, Programming Environment,
Installation ও Syntax পরবর্তী Topic-গুলোতে আলোচনা করা হবে।

\end{itemize}

\end{rememberbox}

\small

%================================================
% 5.10 Fourth Generation Language
%================================================
